
\documentclass[12pt]{article}
\usepackage[english]{babel}
\usepackage[utf8x]{inputenc}
\usepackage{tikz}
\usepackage{tikz-cd}
\usepackage{mathtools}
\usepackage[T1]{fontenc}
\usepackage{listings}
\usepackage{bookmark}


\makeatletter
\def\input@path{{../../style/}}
\makeatother

\usepackage{../../style/quiver}
\makeatletter
\def\input@path{{../../style/}}
\makeatother

\usepackage{../../style/scribe}
\usepackage{fancyhdr}

\usepackage{parskip}
\setlength{\parskip}{1em}
\setlength{\parindent}{0pt}

\DeclareMathOperator{\len}{len}
\newcommand{\fg}{\mathfrak g}
\newcommand{\Rees}{Rees}

\DeclareMathOperator{\Frac}{Frac}
\begin{document}

\lhead{Songyu Ye}
\rhead{\today}
\cfoot{\thepage}

\title{Title}
\author{Songyu Ye}
\date{\today}
\maketitle

\begin{abstract}
    Abstract
\end{abstract}

\tableofcontents

\section{Feb 24}
Let $X$ be complex variety, $X = \bigsqcup X_j$ locally closed decomposition.
\begin{definition}
    A sheaf $\cF$ on $X$ is \emph{constructible} if $\cF|_{X_j}$ is a local system of finite type for each $j$, for some locally closed decomposition $X = \bigsqcup X_j$. Put $D^b_c(X)$ the bounded derived category of constructible sheaves on $X$, those guys with constructible cohomology sheaves.
\end{definition}

\begin{theorem}
    The six functors and Verdier duality preserve $D^b_c(X)$.
\end{theorem}

Special case:
Let $i:X\to Y$ be a locally closed embedding. Recall that we have the pushforward with support $i_!$ defined on stalks by \begin{align*}
(i_!\cF)_y = \begin{cases}
    \cF_x & y = i(x)\\
    0 & y\notin i(X)
\end{cases}
\end{align*}
\begin{proposition}
    $i_!$ is exact.
\end{proposition}
 
\begin{definition}
    The exception pullback $i^!$ is explicitly defined by the formula
\begin{align*}
    U \mapsto \lim_{V \subset Y, V\cap \overline{X}=U} {s\in \cF(V) \mid \mathrm{Supp}(s) \subset V}
\end{align*}
In the special case $i_x:\{x\}\to X$ is the inclusion of a point, we have this computes the costalk
\begin{align*}
i_x^!\cF = \lim_{U\ni x} \{s\in \cF(U) \mid \mathrm{Supp}(s) \subset U\}
\end{align*}
\end{definition}

\begin{proposition}
    Let $i:X\to Y$ be a locally closed embedding. Then there is a derived adjunction
\begin{align*}
    \mathrm{Hom}_{D^b_c(Y)}(i_!\cF, \cG) \cong \mathrm{Hom}_{D^b_c(X)}(\cF, i^!\cG)
\end{align*}
\end{proposition}
Note that if $j$ is a closed embedding $i_! = i_*$, and if $i$ is an open embedding $i^! = i^*$.

\begin{theorem}
    Let $i:Z\to X$ be a closed embedding, $j:U\to X$ the complementary open embedding. Then $i^*\circ j_! = i^! \circ Rj_* = j^* \circ i_* = 0$.
    There are also distinguished triangles
\begin{align*}
    j_!j^! F \to F \to i_*i^* F \to j_!j^! F[1]
\end{align*}
and \begin{align*}
    i_!i^! F \to F \to Rj_*j^* F \to i_!i^! F[1]
\end{align*}
\end{theorem}

\begin{theorem}[Verdier duality]
    Suppose $f:X\to Y$ is locally compact and Hausdorff and $Rf_!$ has finite cohomological dimension. Then there exists a right adjoint $f^!:D^*(Y;k) \to D^*(X;k)$. Let $p:X\to *$ be the collapse map. The dualizing complex is defined by $\omega_X := p^!k$. Then we have a duality functor $\mathbb D: D^*(X;k) \to D^*(X;k)$ defined by $\mathbb D(\cF) = R\mathrm{Hom}(\cF, \omega_X)$.
\end{theorem}

Check. Let $i:Z\to X$ be a closed embedding, $j:U\to X$ the complementary open embedding. Consider $j_! k_{X\setminus Z}$. This is the sheaf of locally constant functions on $X$ supported on $X\setminus Z$. Take further to be $X=S^1$ and $Z=\{p\}$ a point. 

Recall that $j_!$ is defined by the formula for $j$ open
\begin{align*}
    j_!F(U) = \{s\in F(j^{-1}(U)) \mid j\vert_{\mathrm{Supp}(s)}: \mathrm{Supp}(s) \to U \text{ is proper}\}
\end{align*}
On open sets $U \subset X$ not containing $p$ we get locally constant functions. On open sets $U$ containing $p$, we get the locally constant functions, so we get 2 copies of $k$ because removing the point disconnected the interval. Those sections there whose image has proper support in $U$ must be those which vanish everywhere. So we get $j_! k_{X\setminus Z}(U) = 0$ for $U$ containing $p$. 

Check: Give an interpretatino of $H^k(X,Z,k)$ and prove the long exact sequence of relative cohomology.
Claim: $H^*(S^1,*;k) = \Gamma(S^1, j_! k_{S^1\setminus \{p\}})$.

\subsection{}
Consider the nilpotent cone of $\sl(2)$:
\begin{align*}
    \tilde{\mathcal{N}} = \set{ A = \begin{pmatrix}
    x & y\\
    z & -x
\end{pmatrix}, [u:v] \in \mathbb P^1 | Ax = 0} \cong T^*\C\P^1
\end{align*}
Note that specifiying $\ell = [u:v]$ gives a line in $\C^2$, and asking for a nilpotent $A$ to vanish on $\ell$ actually reduces to a map $\C^2/\ell \to \ell$, because Jordan canonical form implies that $A$ nilpotent and traceless implies that $A^2 = 0$ is precisely equivalent to $\im A = \ker A$.

We have the map $\mu: \tilde{\mathcal{N}} \to \mathcal{N}$ the projection to the first factor. \begin{align*}
\mathcal N = \set{ \begin{pmatrix}
    x & y\\
    z & -x
\end{pmatrix} \in \sl(2) \mid x^2 + yz = 0}
\end{align*}
It is the moment map.
What are the fibers of $\tilde{\mathcal{N}}$? Pass to Jordan forms of $N$. The matrix $N = \begin{pmatrix}
    0 & 1\\
    0 & 0
\end{pmatrix}$ has fiber $*$ and the matrix $N = 0$ has fiber $\C\P^1$. 

Consider the complement $U$ of $0$ in $\mathcal N$. Divide out by the scaling action. It will be a real 3 manifold. Reasoning about the (transitive) adjoint action of $SU(2)$ on $U$ we see that $U \cong \R \times L$ with $L = S^3/\Z_2 = \R\P^3 = SO(3)$ is the unit cotangent bundle of $S^2$.

Now consider $R\Gamma k_U$. We observed $(Rj_* k_U)$ is the cohomology of $\R\P^3$, the neighborhood of the cone point (up to homotopy). We can comppute $R\mu_* k_{\tilde{\mathcal{N}}}$ using proper base change. Consider the distinguished triangle
\begin{align*}
    k_{\mathcal N} \to R\mu_* k_{\tilde{\mathcal{N}}} \to i_* k_0[-2]
\end{align*}

Then if $2$ is invertible in $k$, we can split the triangle and get $R\mu_* k_{\tilde{\mathcal{N}}} \cong k_{\mathcal N} \oplus i_* k_0[-2]$. 

Exercise: calculate the costalk $i^!R\mu_* k_{\tilde{\mathcal{N}}}$ and the map $i^! R\mu_* k_{\tilde{\mathcal{N}}} \to i^* R\mu_* k_{\tilde{\mathcal{N}}}$.

We have to parse $i_0^!$. It is sections which vanish relative to the boundary. Then we have the Springer map $\mu$ and we are thinking about pulling back sections. Unwinding, we find that $i_0^! R\mu_* k_{\tilde{\mathcal{N}}}$ identifies with $C^*(DS^2, \partial DS^2)$ is the cohomology of $P^1$ shifted by the Thom class.
\end{document}